\section{再現性と測定物の来歴}
\label{app:reproducibility}

\subsection{交差モデル比較 run の同定情報}
\label{app:benchmark_run}

\cref{sec:results} の交差モデル比較は、
次の 1 回の実行に由来する。

\begin{center}
  \small
  \begin{tabularx}{\linewidth}{@{}p{0.30\linewidth}X@{}}
    \toprule
    項目 & 値 \\
    \midrule
    実行ディレクトリ &
      \texttt{knowledge/runs/run-court-alignment-crossmodel-n120-s42/} \\
    manifest fingerprint &
      \texttt{bfc02859}\allowbreak\texttt{40621d6b6}\allowbreak
      \texttt{aa96fc79b6e308f}\allowbreak\texttt{d24e9d7843}\allowbreak
      \texttt{41f5b7668caa}\allowbreak\texttt{964d77675c} \\
    標本 & domain ごとに決定論的 \resBenchDomainSamples 件、seed \resBenchSeed \\
    推論デバイス & CPU のみ。
      \texttt{CUDA\_VISIBLE\_DEVICES} は空、
      \texttt{torch.cuda.is\_available()} は偽 \\
    実行時間 & \resBenchElapsedSeconds 秒（二つの domain と二つのモデルの合計） \\
    設定 fingerprint &
      \texttt{settings\_fingerprint} として manifest に保存 \\
    生成物 & \texttt{manifest.json}, \texttt{metrics.json},
      \texttt{metrics.csv}, \texttt{metrics.tex},
      \texttt{provenance.json}, \texttt{figures/} \\
    原稿側の数値 & \texttt{results/generated/benchmark\_metrics.tex} \\
    \bottomrule
  \end{tabularx}
\end{center}

\paragraph{domain の定義}
\texttt{real\_validation} は公開実装付属データの held-out validation であり、
公式の test split ではない。
\texttt{synthetic\_test} は \SyntheticDataset の trajectory-group-disjoint test で、
同じ会場群の未学習軌道である。
未知会場ではない。

\paragraph{両モデルの同定情報}
\Method 側の checkpoint は次のとおりである。
checkpoint SHA-256 は
\texttt{b863df1f01f00d2ff00a21d56a461879}\allowbreak
\texttt{e3c5af193a9f32cec47a88d2842f8383}（\resBenchOursCheckpointHashShort で始まる）、
読み込み時は checkpoint の hyper-parameters から
入力解像度と前処理を復元し、
補間は training 側の幾何に合わせる。
\BaselineShort 側は commit \BaselineCommit の実装と
配布重み（SHA-256 \texttt{09aa8c4338459ba1d643f2dc329f45f}\allowbreak
\texttt{464dedec3720fccc1a4abfd1f7b464d04}）を、
補助中心点を除いた十四点で評価した。
どちらも adapter のソースハッシュが
\texttt{provenance.json} に記録されている。

\paragraph{RANSAC の失敗理由を残す}
ホモグラフィが失敗した標本は成功率の分母から除かず、
理由別の件数も \texttt{metrics.json} に保存する。
合成 \BaselineShort の失敗は
有効キーポイント不足 \resSynthTcdInsufficientFailures 件と
RANSAC 失敗 \resSynthTcdRansacFailures 件である。

\paragraph{結果図の同定情報}
\texttt{figures/generated/} の PNG はこの run の
\texttt{figures/} から複製したものである。
SHA-256 は次のとおり。

\begin{center}
  \small
  \begin{tabular}{@{}ll@{}}
    \toprule
    ファイル & SHA-256（先頭 16 桁） \\
    \midrule
    \texttt{summary\_bars.png} & \texttt{9d5d9924f9295d3d} \\
    \texttt{error\_cdf.png} & \texttt{91ec0abeececd7bd} \\
    \texttt{montage\_real\_validation.png} & \texttt{68de4e8ac1623cc9} \\
    \texttt{montage\_synthetic\_test.png} & \texttt{740f32bd04c70ada} \\
    \midrule
    \texttt{classical/real\_outlier\_overlay.png} & \texttt{71de969082dde964} \\
    \texttt{classical/synthetic\_wrong\_fit.png} & \texttt{21e85304a494d630} \\
    \texttt{classical/synthetic\_extreme\_overlay.png} & \texttt{9e56ad70a7ac4185} \\
    \texttt{classical/synthetic\_degenerate\_fit.png} & \texttt{5d42c87bed4b6686} \\
    \bottomrule
  \end{tabular}
\end{center}

montage の列は GT を緑、\Method を橙、\BaselineShort を青で描き、
4 列目で重ねる。
どの標本を描くかは、その domain の対応誤差の分位点から
決定論的に選ぶ。
補助 probe の PNG は代表例を複製したもので、
\cref{fig:classical_probe} が使うのは 3 枚である。
残りは来歴の記録として置く。

\subsection{古典的 baseline 補助 probe の同定情報}
\label{app:classical_probe}

\cref{sec:results_classical} の補助 probe は、
主比較とは別の実行である。

\begin{center}
  \small
  \begin{tabularx}{\linewidth}{@{}p{0.30\linewidth}X@{}}
    \toprule
    項目 & 値 \\
    \midrule
    実行ディレクトリ &
      \texttt{knowledge/runs/run-court-classical-probe-n10/} \\
    upstream & \cite{gchlebusrepo} の commit \resClassicCommitShort \\
    ライセンス & BSD-3-Clause \\
    OpenCV 4 互換 patch & SHA-256
      \texttt{12422aec}\allowbreak\texttt{27f589b9}\allowbreak
      \texttt{2be68971e56f74c}\allowbreak\texttt{a13431decf}\allowbreak
      \texttt{bec7482f0f5cc}\allowbreak\texttt{45a2d8f555} \\
    環境 & \resClassicEnv \\
    実行バイナリ & SHA-256
      \texttt{1a805c96389e6d33}\allowbreak
      \texttt{7a5666b06f38b0d1}\allowbreak
      \texttt{5bf8a73119eb9d3e}\allowbreak
      \texttt{0c4107a136614443} \\
    標本 & \resClassicProbeSamples（主比較とは別） \\
    閾値 & 既定値のまま。検出閾値の調整はしない \\
    権威記録 & \texttt{REPORT.md}, \texttt{sweep\_results.json} \\
    再実行 & \texttt{repro.sh}（作業ディレクトリを \texttt{mktemp} で作る） \\
    \bottomrule
  \end{tabularx}
\end{center}

patch は OpenCV 4 の API 名と CMake の依存解決、
および検証用 overlay 出力の追加だけを行う。
overlay は環境変数を設定したときだけ保存され、
未設定時は upstream と同一の挙動である。
入力は lossless に固定し、
合成入力は正準の \texttt{rgb.f32.npz} から
量子化した画像を使う（元画像との最大絶対差は 0）。

この probe の限界は次のとおりである。
失敗率は 10 件ずつの標本であり、
\texttt{B00} test の \resClassicBzeroTestFrames 枚すべてではない。
幾何の妥当性は「外接矩形がフレームの 10\% 未満なら縮退」という
事前に決めた補助基準と目視のみで、
GT との点対応比較は行っていない。

\subsection{学習 run の同定情報}

\cref{tab:ckpt_validation} の値は次の run に由来する。

\begin{center}
  \small
  \begin{tabular}{@{}ll@{}}
    \toprule
    項目 & 値 \\
    \midrule
    run 記録 & \RunName \\
    実行 commit & \texttt{b8ccc6bac5beac082b97e94c8b42edd9da8d2f40} \\
    作業ブランチ & \texttt{codex/court-colab-multiscale} \\
    設定名 & \texttt{train\_mixed\_multiscale} \\
    checkpoint & \CheckpointName \\
    seed & 42 \\
    記録 step & 29843 \\
    \bottomrule
  \end{tabular}
\end{center}

\paragraph{学習条件}
検出器は \BackboneName を凍結して入力に使い、
タスク Transformer と DPT 系デコーダ、
深さ 3 の残余密ヘッドを用いる。
入力は長辺 256、320、384、448、512 を混ぜた multiscale である。
再開元は同じ出力ディレクトリの version 1 の最終 checkpoint である。

\paragraph{混合条件}
学習は実写と合成の混合であり、
バッチ内の固定比率は合成 4、実写 4 である。
混合の指定は設定ファイル
\texttt{train\_mixed\_multiscale.yaml} と run 記録にあり、
runner は学習前に混合節を分離してから
標準側の設定だけを保存する。
そのため checkpoint と一緒に保存される
\texttt{hparams.yaml} には混合節が残らず、
このファイルだけを見ると合成のみの設定に見える。
したがって本稿は、この run を
\emph{合成のみで学習した run} とは記述しない。
学習データは実写と合成の両方を含む。

\paragraph{データ}
合成側は \SyntheticDataset の \SceneSet で、
採用フレームは trajectory-group-disjoint に分割されている
（\cref{tab:dataset_summary}）。
実写側は公開実装付属の train と validation を用いる。
独立した test split は合成側にのみ存在し、
実写側には独立した test が無い。
公開実写データは train \TCDTrainCount / validation \TCDValCount である。
変換時に、十四点注釈が縮退してコート平面を定められない
1 サンプルを学習から除外する。

\subsection{run の停止}

この run は epoch 18 の 91\% で失敗した。
データ読み込み時に、
\texttt{npz} アーカイブ内の \texttt{byte\_planes.npy} で
CRC-32 の不一致が発生し、
\texttt{BadZipFile} として送出された。

したがって次の二点を明示する。

\begin{enumerate}[leftmargin=1.6em]
  \item この run は設定上の最大 epoch まで完走していない。
  \item \cref{tab:ckpt_validation} の値は、
  epoch \resCkptEpoch の checkpoint 選択時に記録された
  synthetic validation であり、
  完走した最良モデルの成績ではない。
\end{enumerate}

今後の比較では、
同じ停止を避けるためにデータの健全性検査を学習前に行い、
停止した run は
結果を見た後に取り除かず、
そのまま記録する。

\subsection{ベースラインの同定情報}

\begin{table}[h]
  \centering
  \caption{\textbf{比較に用いる公開実装。}
  コードと重みは再配布せず、
  固定した commit と取得物のハッシュのみを記録する。
  }
  \label{tab:baseline_ids}
  \small
  \begin{tabularx}{\linewidth}{@{}p{0.24\linewidth}X@{}}
    \toprule
    項目 & 値 \\
    \midrule
    リポジトリ & \BaselineName\cite{tcdrepo} \\
    commit & \texttt{e5cd4f1ce26b15361700d3d89e068cbf0e82749e} \\
    構造 & 補助中心点を含む十五チャネルのヒートマップ回帰 \\
    入力解像度 & 640\(\times\)360 \\
    付属データ & \TCDImageCount 枚、十四点注釈、
      train \TCDTrainCount / validation \TCDValCount \\
    ライセンス & リポジトリに確認できるライセンス表記が無い。
      論文も確認できない \\
    本研究での扱い & 実装ベースライン。コードと重みは再配布しない \\
    公式重み & 配布物から取得し SHA-256
      \texttt{09aa8c4338459ba1d643f2dc329f45f}\allowbreak
      \texttt{464dedec3720fccc1a4abfd1f7b464d04} \\
    \midrule
    補助 probe & 古典的コートモデル当てはめ\cite{farin2004courtmodels}の
      派生実装\cite{gchlebusrepo} \\
    commit & \texttt{762d077541a77abf4923f5f8f689a1410927d35e} \\
    ライセンス & BSD-3-Clause \\
    状態 & CPU のみで実行済み。主比較とは別標本の補助 probe
      （付録~\ref{app:classical_probe}） \\
    \bottomrule
  \end{tabularx}
\end{table}

\paragraph{その他の比較候補}
サッカー向けの校正手法は、
テンプレートと出力の形式が異なるため、
そのままテニスの定量比較へ混ぜない。
比較する場合は、
対象のプレー面と出力の型を揃えた条件を別に定義する。

\subsection{ビルドと再現}

本原稿は LuaLaTeX と luatexja で組版する。
数値は \texttt{results/generated/} に置いた
\texttt{benchmark\_metrics.tex} が定義するマクロだけを参照し、
本文・表・図へ直接書き込まない。
したがって数値を更新する場合は、
この 1 ファイルを差し替えれば本文と表が追従する。

\texttt{make} は本文を 1 回ビルドし、
\texttt{make check} は未定義の参照・引用、
TeX エラー、欠落した入力ファイル、
40\,pt を超える Overfull box を検査する。
\texttt{make watch} は \texttt{latexmk -pvc} で監視し、
TeX ソースだけでなく
\texttt{figures/generated/} の PNG が変わったときにも再ビルドする。
\texttt{Makefile} の \texttt{SOURCES} が
生成図の PNG を wildcard で含むためである。
ビルドに成功した時点で
\texttt{build/main.pdf} がトップレベルの \texttt{main.pdf} へ複製され、
失敗時は複製しない。

\paragraph{測定済みであることの担保}
\texttt{resultsmeasured} と \texttt{resultsfigures} は
この版で真にしてある。
未測定として残るのは、
\cref{tab:ablation_plan} の視点分布 ablation と
採用視点分布の集計図だけであり、
これらは「未測定」または「未作成」と本文で明示する。
値があるように見せる既定値は置かない。
